\documentclass[12pt, a4paper]{letter}
\usepackage[utf8]{inputenc}
\usepackage[T1]{fontenc}
\usepackage{geometry}
\usepackage{amsmath, amssymb}
\usepackage{hyperref}

\geometry{margin=1in}

\hypersetup{colorlinks=true, linkcolor=blue, urlcolor=blue}

\signature{Rafael D. De Paz \\ Independent Researcher \\ \texttt{me@rdepaz.com}}
\address{Rafael D. De Paz \\ Independent Researcher \\ \url{https://logos.pub}}

\begin{document}
\begin{letter}{Editorial Board \\ Annals of Mathematics \\ Princeton University \& IAS}

\opening{Dear Editors,}

I am writing to submit the enclosed manuscript, entitled

\begin{center}
\textbf{``The Riemann Hypothesis via Discrete Spectral Theory: \\
Zeros of the Zeta Function as Eigenvalues of a Self-Adjoint Lattice Operator,''}
\end{center}

\noindent for consideration for publication in \textit{Annals of Mathematics}.

\medskip

\textbf{Summary.} The paper pursues the Hilbert-P\'olya approach to the
Riemann Hypothesis: constructing an explicit self-adjoint operator whose
spectrum corresponds to the non-trivial zeros of the Riemann zeta function
$\zeta(s)$. We formulate the problem on a discrete spectral lattice,
where self-adjointness automatically forces all eigenvalues to be real---
placing all zeros on the critical line $\mathrm{Re}(s) = 1/2$.

\medskip

\textbf{Approach.} The construction is motivated by: (i) Montgomery's pair
correlation conjecture, which links zero spacings to GUE random matrix
eigenvalue statistics; (ii) the work of Berry and Keating on the classical
Hamiltonian $H = xp$; and (iii) the spectral geometry of arithmetic surfaces.
The discrete formulation provides an explicit, finite-dimensional Hamiltonian
whose adjacency spectrum converges to the zeta zeros as the lattice refines.

\medskip

\textbf{Suggested Reviewers.}
\begin{enumerate}
  \item \textbf{Prof. Alain Connes} (IHES/Coll\`ege de France) --- Author of
    the spectral interpretation of the Riemann zeros via noncommutative geometry.
  \item \textbf{Prof. Peter Sarnak} (Princeton/IAS) --- Leading expert on
    $L$-functions, random matrix theory, and the Riemann Hypothesis.
  \item \textbf{Prof. Michael Berry} (Bristol) --- Developer of the Berry-Keating
    Hamiltonian approach to the Riemann Hypothesis.
\end{enumerate}

\medskip

This manuscript has not been previously published and is not under consideration
at any other journal. I confirm no conflicts of interest.

\closing{Respectfully submitted,}
\end{letter}
\end{document}
